\documentclass[11pt]{article}
\usepackage[margin=1in]{geometry}
\usepackage{amsmath,amssymb}
\usepackage{enumitem}
\usepackage{booktabs}
\usepackage{framed}

\newcommand{\indep}{{\bot\negthickspace\negthickspace\bot}}
\DeclareMathOperator{\E}{\mathbb{E}}

\pagestyle{plain}

\begin{document}

\begin{center}
{\Large \textbf{PS200C: Quantitative Methods}} \\[6pt]
{\large In-class Activity 8: Instrumental Variables}
\end{center}

\vspace{4pt}
\noindent\textbf{Name:} \hrulefill

\vspace{12pt}

%%% SETUP for Q1--Q2 %%%
\begin{framed}
\noindent\textbf{Setup for Q1--Q2.}  Suppose a city is studying whether being incarcerated (vs.\ a non-custodial sentence) for a first offense affects re-arrest within five years.  In this jurisdiction there are two judges who hear first-offense cases: \textbf{Judge A} (relatively lenient) and \textbf{Judge B} (relatively stringent).  Defendants are assigned (essentially at random) to one of the two.

\medskip
The researcher proposes:
\begin{itemize}[itemsep=0pt, topsep=2pt]
\item $Z = 1$ if assigned to Judge B (stringent); $Z = 0$ if assigned to Judge A (lenient).
\item $D = 1$ if incarcerated; $D = 0$ if non-custodial sentence.
\item $Y = 1$ if re-arrested within five years.
\end{itemize}
\end{framed}

\vspace{6pt}

%%% QUESTION 1 %%%
\noindent\textbf{Question 1: Who are the compliers?}  In one sentence, describe the compliers in this design --- the subgroup whose treatment status would be moved by the instrument.

\vspace{4cm}

%%% QUESTION 2 %%%
\noindent\textbf{Question 2: State the exclusion restriction.}  In one sentence, write what the exclusion restriction \emph{requires} here.  Can you think of anything that might violate it in this setting?

\vspace{6cm}

\newpage

%%% QUESTION 3 %%%
\noindent\textbf{Question 3: Understanding the assumptions.}  A researcher proposes using \textit{distance to the nearest college} as an instrument for years of education, with adult earnings as the outcome.

\begin{itemize}
\item How might the ignorability assumption be violated here?
\vspace{2.5cm}
\item How might the exclusion restriction be violated here?
\vspace{2.5cm}
\end{itemize}

%%% QUESTION 4 %%%
\noindent\textbf{Question 4: Wald by hand.}  Suppose you have a randomized encouragement-design experiment.  $Z=1$ if assigned to receive a free flu-shot voucher; $D=1$ if the person actually got vaccinated; $Y$ is days of missed work that flu season.

\medskip
\begin{center}
\begin{tabular}{lcc}
\toprule
            & \textbf{$Z=0$ (no voucher)} & \textbf{$Z=1$ (voucher)} \\
\midrule
$\bar D$    &   0.20    &   0.55    \\
$\bar Y$    &   3.0 days  &  2.3 days \\
\bottomrule
\end{tabular}
\end{center}
\medskip

\begin{enumerate}[label=(\alph*), itemsep=0.4em]
\item Compute the ITT (effect of $Z$ on $Y$) and the first stage (effect of $Z$ on $D$).

\vspace{2.5cm}

\item Compute the Wald / IV estimate for the effect of vaccination on missed work.

\vspace{2.5cm}

\item For whom is this the treatment effect? (Describe in a sentence)

\vspace{3cm}
\end{enumerate}

%\newpage
% %%% QUESTION 5 %%%
% \noindent\textbf{Question 5: Why doesn't IV identify the ATE?}  In a paragraph (3--5 sentences), explain why the Wald estimator recovers the LATE on compliers and not the average treatment effect on the whole population.  Make sure to use the words \textit{monotonicity}, \textit{compliers}, and \textit{always-takers / never-takers} somewhere in your answer.

% \vspace{10cm}

\end{document}
